\section{Отчёт реммастера}

Ниже приведён список поломок и методов ремонта, применявшихся в походе:

\begin{itemize}
	\item \textbf{Лыжное снаряжение:}
	\begin{itemize}
		\item Натяжитель тросика («лягушка») – потеряна штатная гайка крепления тросика, т.к. на винте не было фиксатора и, видимо, при пешем переходе гайка скрутилась с винта. Отремонтирована с помощью шайбы толщиной 1.5~мм, у которой загнули края, чтобы она фиксировала тросик. Саму шайбу фиксировали гайкой М5. И гайка, и шайба нашлись в избе Петра Севостьянова, с ремонтом помогали обитатели избы, за что им большое спасибо.
		\item Лыжное крепление – 1 шт. (отодрано от платформы, прикручено заново на биваке)
		\item Рвались тросики лыжных креплений, оперативно заменялись за запасные, под длине были подобраны заранее
		\item В первый день, в рамках подгонки снаряжения, для лучшего натяжения на некоторых лыжах были переставлены лягушки
		\item У многих участников была проблема отклеивания камуса от поверхности лыжи. Оказалось, что армированный скотч , двусторониий скотч и цианоакрилатный клей плохо подходят для решения этой проблемы (на морозе не клеят). Лучше всего подходил обычный скотч, которым прихватывались боковые поверхности лыжи вместе с камусами. Такой ремонт держался не более одного ходового дня, так как металлический кант легко рвал скотч, но в целом, проблему это решало.
	\end{itemize}
	
	\item \textbf{Личное снаряжение и одежда:}
	\begin{itemize}
		\item Спальник (прожгли). Ремонтировали армированным скотчем (скорее безуспешно), специальными заплатками (несколько успешнее)
		\item Пуховик (два раза рвался, заклеивали заплатками)
	\end{itemize}
\end{itemize}

По поводу сверления лыж сделали следующие наблюдения:

\begin{itemize}

		\item Пластик лыжи легко сверлится шилом
		\item Металлическая скитурная плита шилом не сверлится, что ограничивает возможность крепления камусов на саморезы в полевых условиях

\end{itemize}




\clearpage